\title{Mixed States and Density Matrices}
\author{Old notes}
\date{today}

\begin{document}
\maketitle

\section{Probabilistic Mixtures of Pure Quantum States}

Previously, we encountered only pure states $\ket{\psi}$, described by unit vectors in $\mathbb{C}^d$. Consider the pure state $\ket{+}$:
\begin{equation}
H\ket{0} = \frac{1}{\sqrt{2}} \begin{pmatrix} 1 \ 1 \end{pmatrix} = \frac{1}{\sqrt{2}}(\ket{0}+\ket{1}) = \ket{+}.
\end{equation}

Measuring $\ket{+}$ in the computational basis yields a random outcome, whereas measuring it in the diagonal basis ${\ket{+},\ket{-}}$ yields a deterministic outcome.

Not all quantum systems can be described adequately by pure states. In particular, we wish to describe probabilistic quantum states whose measurement outcomes remain random in all bases. This motivates the introduction of mixed states.

\subsection{Mixed States}

Mixed states are probabilistic mixtures of pure quantum states. They arise naturally in the presence of noise and are essential for describing subsystems of larger quantum systems, particularly entangled systems. Mixed states provide a formalism to describe individual subsystems that cannot be characterized by pure states alone.

\subsection{Probabilistic Mixtures}

Suppose a qubit is assigned state $\ket{0}$ or $\ket{+}$ depending on the outcome of an unknown coin toss. Our prior belief is represented by the distribution
\begin{equation}
\left(\tfrac{1}{2}, \ket{0}\right), \quad \left(\tfrac{1}{2}, \ket{+}\right).
\end{equation}



For complex systems, such descriptions quickly become unwieldy, motivating a more compact representation.

\section{Density Matrices}

\subsection{Background}

\paragraph{Spectral Theorem.} Let $A \in \mathbb{C}^{d \times d}$ be Hermitian. Then there exists an orthonormal basis ${\ket{b_1},\ldots,\ket{b_d}}$ and real eigenvalues $\lambda_1,\ldots,\lambda_d$ such that
\begin{equation}
  A = \sum_j \lambda_j \ket{b_j}\bra{b_j}.
\end{equation}

\paragraph{Positive Semidefinite.} A Hermitian matrix $A$ is positive semidefinite (PSD) if all eigenvalues are nonnegative.

\subsection{Definition}

A density matrix is a matrix $\rho \in \mathbb{C}^{d \times d}$ such that $\rho$ is PSD and $\Tr(\rho)=1$.

A pure state $\ket{\psi}$ corresponds to the density matrix $\rho = \ket{\psi}\bra{\psi}$.

More generally, a probabilistic mixture ${(p_i,\ket{\psi_i})}$ corresponds to
\begin{equation}
  \rho = \sum_i p_i \ket{\psi_i}\bra{\psi_i}, \qquad \sum_i p_i = 1.
\end{equation}

\subsection{Examples}

\paragraph{Example 1.} Density matrices for $\ket{0}$ and $\ket{1}$:
\begin{equation}
  \ket{0}\bra{0} = \begin{pmatrix} 1 & 0 \ 0 & 0 \end{pmatrix}, \quad
  \ket{1}\bra{1} = \begin{pmatrix} 0 & 0 \ 0 & 1 \end{pmatrix}.
\end{equation}

\paragraph{Example 2.} Mixture $\left(\tfrac12,\ket{0}\right)$ and $\left(\tfrac12,\ket{+}\right)$:
\begin{equation}
  \rho = \frac12 \ket{0}\bra{0} + \frac12 \ket{+}\bra{+} = \begin{pmatrix} \tfrac34 & \tfrac14 \ \tfrac14 & \tfrac14 \end{pmatrix}.
\end{equation}

\paragraph{Example 3.} Distinct mixtures can yield the same density matrix:
\begin{equation}
  \frac12 \ket{0}\bra{0} + \frac12 \ket{1}\bra{1} = \frac{I}{2} = \frac12 \ket{+}\bra{+} + \frac12 \ket{-}\bra{-}.
\end{equation}

\subsection{Mixtures of Density Matrices}

Given density matrices $\rho$ and $\sigma$ with probabilities $p$ and $q$, the mixture
\begin{equation}
  \tau = p\rho + q\sigma
\end{equation}

is again a valid density matrix.

\subsection{Projective Measurements}

A projective measurement is a set ${M_1,\ldots,M_k}$ of Hermitian projection operators satisfying
\begin{equation}
  M_i^2 = M_i, \quad M_i^\dagger = M_i, \quad \sum_i M_i = I.
\end{equation}

Measuring $\rho$ yields outcome $i$ with probability $\Tr(M_i \rho)$.

\subsection{Unitary Evolution}

Under unitary evolution $U$, a density matrix evolves as
\begin{equation}
  \rho \mapsto U \rho U^\dagger.
\end{equation}

\subsection{Observables}

Observables are Hermitian operators $A$. The expectation value of $A$ in state $\rho$ is
\begin{equation}
  \Tr(A\rho) = \sum_j \lambda_j \Pr(\text{outcome } \ket{b_j}).
\end{equation}

\subsection{Multiple Systems}

The joint state of systems with density matrices $\rho$ and $\sigma$ is $\rho \otimes \sigma$. Not all composite states are tensor products, but such states need not be entangled.

\section{Traces and Partial Traces}

Given a bipartite system $AB$ with density matrix $\rho_{AB}$, the reduced state on $A$ is obtained by the partial trace
\begin{equation}
  \rho_A = \Tr_B(\rho_{AB}).
\end{equation}

Every mixed state can be viewed as the partial trace of a pure state on a larger system (a purification).

\subsection{Distinguishability of Density Matrices}

The trace distance between $\rho$ and $\sigma$ is defined as
\begin{equation}
  D(\rho,\sigma) = \frac12 | \rho - \sigma |_1 = \frac12 \Tr |\rho-\sigma|.
\end{equation}

The trace distance quantifies the optimal distinguishability of two quantum states and satisfies symmetry, the triangle inequality, monotonicity under partial trace, and unitary invariance.

\end{document}
